\question[5]
Beim asymmetrischen Verfahren muss sichergestellt sein, dass ein öffentlicher Schlüssel wirklich zu
der Person (oder dem Server) gehört, die ihn präsentiert. Dazu dient eine
\textit{Public-Key-Infrastruktur} (PKI) mit Zertifikaten, die von einem Trust-Center ausgestellt werden. \\
a. Erläutere, welches Problem ein Zertifikat löst und wie ein Trust-Center dabei vorgeht. \\
b. Nenne zwei Anwendungsbereiche von Public-Key-Infrastrukturen aus dem Alltag. \\
c. Bewerte einen Sicherheitsaspekt: Was passiert, wenn der private Schlüssel eines Trust-Centers
in falsche Hände gerät?

\begin{solutionbox}{6cm}
a. Ein Zertifikat bestätigt die \textit{Authentizität} eines öffentlichen Schlüssels, also die
Zuordnung Schlüssel $\leftrightarrow$ Inhaber. Das Trust-Center prüft die Identität des Antragstellers
und signiert dann dessen öffentlichen Schlüssel zusammen mit den Inhaberdaten mit seinem eigenen
privaten Schlüssel. Wer dem Trust-Center vertraut und dessen öffentlichen Schlüssel kennt, kann
das Zertifikat prüfen. \\
b. Zum Beispiel: Transportverschlüsselung im Web (HTTPS / TLS-Zertifikate), Mailverschlüsselung
(S/MIME), digitaler Personalausweis, Signaturkarten. \\
c. Dann kann der Angreifer beliebige gefälschte Zertifikate ausstellen und sich als jede beliebige
Stelle ausgeben (Man-in-the-Middle). Das Vertrauen in alle von diesem Trust-Center ausgestellten
Zertifikate ist zerstört; sie müssen zurückgezogen werden.
\end{solutionbox}
